\begin{table}[!htbp]
\small
\setlength{\tabcolsep}{5pt}

\caption{\label{tab:income-definition-robustness}Robustness to income definition: inflation gap (in points, 2022)}
\centering
\begin{tabular}[t]{lrrr}
\toprule
Country & Total household income & Equivalised income & Difference\\
\midrule
Austria & 0.2 & 0.2 & 0.1\\
Belgium & 3.2 & 1.4 & -1.8\\
Croatia & 0.3 & 0.8 & 0.5\\
Cyprus & 1.0 & 0.6 & -0.4\\
Estonia & 5.0 & 5.4 & 0.5\\
\addlinespace
Finland & -0.6 & -0.4 & 0.3\\
France & 0.3 & 0.1 & -0.2\\
Germany & 1.8 & 1.7 & 0.0\\
Greece & 1.4 & 0.7 & -0.7\\
Ireland & 4.0 & 2.1 & -1.8\\
\addlinespace
Latvia & 5.2 & 4.0 & -1.2\\
Lithuania & 4.6 & 3.3 & -1.4\\
Luxembourg & -0.1 & 0.0 & 0.1\\
Malta & 0.2 & 0.0 & -0.2\\
Netherlands & 3.9 & 2.6 & -1.3\\
\addlinespace
Slovakia & 0.1 & 0.2 & 0.1\\
Slovenia & 0.6 & 0.4 & -0.2\\
Spain & 1.2 & 0.8 & -0.4\\
Euro area (18 countries) & 1.4 & 1.1 & -0.3\\
\bottomrule
\end{tabular}
\begin{flushleft}
\footnotesize{\emph{Notes:} The table reports first-quintile minus fifth-quintile annual-average household inflation in percentage points. Both columns use the same 2022 household inflation observations, the same rent-excluded consumption universe, RAS weights, survey weights, and country samples. Only the national quintile ranking changes. Total household income is EUR\_HH099. Equivalised income divides EUR\_HH099 by the OECD equivalence scale HB061. Difference is equivalised income minus total household income. The euro-area row averages national gaps with 2022 HICP country weights over the 18-country sample.}
\end{flushleft}
\end{table}
